% ID: SEL-D0014-BA-E02
% ESTADO: cotejada. Segunda pasada visual completa y borrador compilado revisados.
% Fuente: Selectividad/matematicas_aplicadas/matematicas_aplicadas_P_B.pdf
% SHA-256: 51601c4f5a6c1dfcd860f621a659f8b9c09f53ff4cd31a1fa1e827dbd8f06159
% Página: 1; bloque A; ejercicio original 2.

Se consideran las matrices
$A=\begin{pmatrix}2&1&-1\\a&-1&-1\\3&0&-2a\end{pmatrix}$ y
$B=\begin{pmatrix}1&-1\\2&0\\1&-2\end{pmatrix}$.
\begin{enumerate}[label=\alph*)]
  \item \textbf{(0.7 puntos)} Determine para qué valores de $a$ tiene inversa
  la matriz $A$.
  \item \textbf{(1 punto)} Para $a=2$, calcule la matriz inversa de $A$.
  \item \textbf{(0.8 puntos)} Para $a=0$, resuelva la ecuación matricial
  $X\cdot A^{-1}-B\cdot B^t=I_3$.
\end{enumerate}
